\documentclass[reqno]{amsart}
\usepackage[T1]{fontenc}
\usepackage{amsmath,amssymb,amsthm}
\usepackage{xr-hyper}
\usepackage[hidelinks]{hyperref}
\externaldocument[ds-]{20261125T004000--denominators-and-sha__workbook}
\externaldocument[cg-]{20261125T014100--congruence-modules__workbook}
\externaldocument[tt-]{20261125T004100--tate-twists-and-sha__workbook}
\externaldocument[gr-]{20261125T011000--galois-representations-black-box__workbook}
\externaldocument[tf-]{20261126T031000--twists-with-few-congruences__workbook}
\externaldocument[oc-]{20261126T020430--the-minus-part-by-gauss-sums__workbook}
\externaldocument[c32-]{20261126T030300--the-data__workbook}

% Chapter number; \ref{chap} gives it.
\newcounter{chap}
% Chapter numbers, generated from 20261125T012000--grh-workbook__guide.org
% by scripts/make-chapters.py.  Do not edit: change the order of the
% index and rerun the script (./preview-site.sh and CI do this).
%
% Usage in a sheet, after \newcounter{chap}:
%   % Chapter numbers, generated from 20261125T012000--grh-workbook__guide.org
% by scripts/make-chapters.py.  Do not edit: change the order of the
% index and rerun the script (./preview-site.sh and CI do this).
%
% Usage in a sheet, after \newcounter{chap}:
%   % Chapter numbers, generated from 20261125T012000--grh-workbook__guide.org
% by scripts/make-chapters.py.  Do not edit: change the order of the
% index and rerun the script (./preview-site.sh and CI do this).
%
% Usage in a sheet, after \newcounter{chap}:
%   \input{chapters}\chapnumber{<sheet ID>}
\makeatletter
\providecommand\chapters@set[2]{\@namedef{chapters@#1}{#2}}
\chapters@set{20261125T001000}{1}% zeros-to-ramp-sums
\chapters@set{20261125T001100}{2}% inert-primes-and-estimate-b
\chapters@set{20261125T005000}{3}% galois-fourier-and-averaging
\chapters@set{20261125T005500}{4}% columns-from-a-sphere
\chapters@set{20261125T000000}{5}% slope-method
\chapters@set{20261126T030100}{6}% a-twist-for-each-prime
\chapters@set{20261125T003000}{7}% the-chain-and-the-square-root
\chapters@set{20261125T010000}{8}% rank-lemma-and-zero-estimate
\chapters@set{20261125T010500}{9}% poincare-series-and-near-orthogonality
\chapters@set{20261124T223000}{10}% elliptic-curves-for-the-proof
\chapters@set{20261125T009500}{11}% fock-model-and-theta-intertwiners
\chapters@set{20261124T223100}{12}% complex-multiplication-and-inert-reduction
\chapters@set{20261125T008000}{13}% one-cm-determinant
\chapters@set{20261124T214500}{14}% supersingular-theta-divisibility
\chapters@set{20261124T214600}{15}% the-isotropic-plane
\chapters@set{20261125T013700}{16}% pairings-on-a-formal-group-of-height-two
\chapters@set{20261125T013500}{17}% hecke-nondegeneracy-on-the-tree
\chapters@set{20261126T030200}{18}% the-bound-is-attained
\chapters@set{20261125T009000}{19}% kronecker-limit-and-chowla-selberg
\chapters@set{20261125T002000}{20}% definite-theta-lifts
\chapters@set{20261125T015500}{21}% siegel-weil-in-the-smallest-case
\chapters@set{20261125T002100}{22}% definite-theta-lifts-and-norms
\chapters@set{20261125T016500}{23}% hecke-operators-and-theta-lifts
\chapters@set{20261125T016700}{24}% divisibility-along-the-chain
\chapters@set{20261125T004000}{25}% denominators-and-sha
\chapters@set{20261125T004100}{26}% tate-twists-and-sha
\chapters@set{20261126T020430}{27}% the-minus-part-by-gauss-sums
\chapters@set{20261125T014100}{28}% congruence-modules
\chapters@set{20261125T014500}{29}% eisenstein-congruences-and-selmer-classes
\chapters@set{20261125T014700}{30}% pseudocharacters-and-the-length-bound
\chapters@set{20261125T011000}{31}% galois-representations-black-box
\chapters@set{20261126T030300}{32}% the-data
\chapters@set{20261126T020210}{33}% bessel-poincare-series-and-their-gram-operators
\chapters@set{20261126T030400}{34}% the-seesaw-identity
\chapters@set{20261126T030500}{35}% divisibility-of-the-composite-lift
\chapters@set{20261126T030600}{36}% theorem-a-short-forms
\chapters@set{20261126T030700}{37}% the-determinant-of-a-chain
\chapters@set{20261126T030800}{38}% the-cost-of-one-lift
\chapters@set{20261126T030900}{39}% denominators-of-projections-to-lifts
\chapters@set{20261126T031000}{40}% twists-with-few-congruences
\chapters@set{20261126T031100}{41}% theorem-c-and-the-main-inequality
\providecommand\chapnumber[1]{%
  \@ifundefined{chapters@#1}%
    {\PackageError{chapters}{Sheet #1 is not listed in the index}%
      {List it in 20261125T012000--grh-workbook__guide.org\MessageBreak
       and run scripts/make-chapters.py.}}%
    {\edef\thechap{\@nameuse{chapters@#1}}}}
\makeatother
\chapnumber{<sheet ID>}
\makeatletter
\providecommand\chapters@set[2]{\@namedef{chapters@#1}{#2}}
\chapters@set{20261125T001000}{1}% zeros-to-ramp-sums
\chapters@set{20261125T001100}{2}% inert-primes-and-estimate-b
\chapters@set{20261125T005000}{3}% galois-fourier-and-averaging
\chapters@set{20261125T005500}{4}% columns-from-a-sphere
\chapters@set{20261125T000000}{5}% slope-method
\chapters@set{20261126T030100}{6}% a-twist-for-each-prime
\chapters@set{20261125T003000}{7}% the-chain-and-the-square-root
\chapters@set{20261125T010000}{8}% rank-lemma-and-zero-estimate
\chapters@set{20261125T010500}{9}% poincare-series-and-near-orthogonality
\chapters@set{20261124T223000}{10}% elliptic-curves-for-the-proof
\chapters@set{20261125T009500}{11}% fock-model-and-theta-intertwiners
\chapters@set{20261124T223100}{12}% complex-multiplication-and-inert-reduction
\chapters@set{20261125T008000}{13}% one-cm-determinant
\chapters@set{20261124T214500}{14}% supersingular-theta-divisibility
\chapters@set{20261124T214600}{15}% the-isotropic-plane
\chapters@set{20261125T013700}{16}% pairings-on-a-formal-group-of-height-two
\chapters@set{20261125T013500}{17}% hecke-nondegeneracy-on-the-tree
\chapters@set{20261126T030200}{18}% the-bound-is-attained
\chapters@set{20261125T009000}{19}% kronecker-limit-and-chowla-selberg
\chapters@set{20261125T002000}{20}% definite-theta-lifts
\chapters@set{20261125T015500}{21}% siegel-weil-in-the-smallest-case
\chapters@set{20261125T002100}{22}% definite-theta-lifts-and-norms
\chapters@set{20261125T016500}{23}% hecke-operators-and-theta-lifts
\chapters@set{20261125T016700}{24}% divisibility-along-the-chain
\chapters@set{20261125T004000}{25}% denominators-and-sha
\chapters@set{20261125T004100}{26}% tate-twists-and-sha
\chapters@set{20261126T020430}{27}% the-minus-part-by-gauss-sums
\chapters@set{20261125T014100}{28}% congruence-modules
\chapters@set{20261125T014500}{29}% eisenstein-congruences-and-selmer-classes
\chapters@set{20261125T014700}{30}% pseudocharacters-and-the-length-bound
\chapters@set{20261125T011000}{31}% galois-representations-black-box
\chapters@set{20261126T030300}{32}% the-data
\chapters@set{20261126T020210}{33}% bessel-poincare-series-and-their-gram-operators
\chapters@set{20261126T030400}{34}% the-seesaw-identity
\chapters@set{20261126T030500}{35}% divisibility-of-the-composite-lift
\chapters@set{20261126T030600}{36}% theorem-a-short-forms
\chapters@set{20261126T030700}{37}% the-determinant-of-a-chain
\chapters@set{20261126T030800}{38}% the-cost-of-one-lift
\chapters@set{20261126T030900}{39}% denominators-of-projections-to-lifts
\chapters@set{20261126T031000}{40}% twists-with-few-congruences
\chapters@set{20261126T031100}{41}% theorem-c-and-the-main-inequality
\providecommand\chapnumber[1]{%
  \@ifundefined{chapters@#1}%
    {\PackageError{chapters}{Sheet #1 is not listed in the index}%
      {List it in 20261125T012000--grh-workbook__guide.org\MessageBreak
       and run scripts/make-chapters.py.}}%
    {\edef\thechap{\@nameuse{chapters@#1}}}}
\makeatother
\chapnumber{<sheet ID>}
\makeatletter
\providecommand\chapters@set[2]{\@namedef{chapters@#1}{#2}}
\chapters@set{20261125T001000}{1}% zeros-to-ramp-sums
\chapters@set{20261125T001100}{2}% inert-primes-and-estimate-b
\chapters@set{20261125T005000}{3}% galois-fourier-and-averaging
\chapters@set{20261125T005500}{4}% columns-from-a-sphere
\chapters@set{20261125T000000}{5}% slope-method
\chapters@set{20261126T030100}{6}% a-twist-for-each-prime
\chapters@set{20261125T003000}{7}% the-chain-and-the-square-root
\chapters@set{20261125T010000}{8}% rank-lemma-and-zero-estimate
\chapters@set{20261125T010500}{9}% poincare-series-and-near-orthogonality
\chapters@set{20261124T223000}{10}% elliptic-curves-for-the-proof
\chapters@set{20261125T009500}{11}% fock-model-and-theta-intertwiners
\chapters@set{20261124T223100}{12}% complex-multiplication-and-inert-reduction
\chapters@set{20261125T008000}{13}% one-cm-determinant
\chapters@set{20261124T214500}{14}% supersingular-theta-divisibility
\chapters@set{20261124T214600}{15}% the-isotropic-plane
\chapters@set{20261125T013700}{16}% pairings-on-a-formal-group-of-height-two
\chapters@set{20261125T013500}{17}% hecke-nondegeneracy-on-the-tree
\chapters@set{20261126T030200}{18}% the-bound-is-attained
\chapters@set{20261125T009000}{19}% kronecker-limit-and-chowla-selberg
\chapters@set{20261125T002000}{20}% definite-theta-lifts
\chapters@set{20261125T015500}{21}% siegel-weil-in-the-smallest-case
\chapters@set{20261125T002100}{22}% definite-theta-lifts-and-norms
\chapters@set{20261125T016500}{23}% hecke-operators-and-theta-lifts
\chapters@set{20261125T016700}{24}% divisibility-along-the-chain
\chapters@set{20261125T004000}{25}% denominators-and-sha
\chapters@set{20261125T004100}{26}% tate-twists-and-sha
\chapters@set{20261126T020430}{27}% the-minus-part-by-gauss-sums
\chapters@set{20261125T014100}{28}% congruence-modules
\chapters@set{20261125T014500}{29}% eisenstein-congruences-and-selmer-classes
\chapters@set{20261125T014700}{30}% pseudocharacters-and-the-length-bound
\chapters@set{20261125T011000}{31}% galois-representations-black-box
\chapters@set{20261126T030300}{32}% the-data
\chapters@set{20261126T020210}{33}% bessel-poincare-series-and-their-gram-operators
\chapters@set{20261126T030400}{34}% the-seesaw-identity
\chapters@set{20261126T030500}{35}% divisibility-of-the-composite-lift
\chapters@set{20261126T030600}{36}% theorem-a-short-forms
\chapters@set{20261126T030700}{37}% the-determinant-of-a-chain
\chapters@set{20261126T030800}{38}% the-cost-of-one-lift
\chapters@set{20261126T030900}{39}% denominators-of-projections-to-lifts
\chapters@set{20261126T031000}{40}% twists-with-few-congruences
\chapters@set{20261126T031100}{41}% theorem-c-and-the-main-inequality
\providecommand\chapnumber[1]{%
  \@ifundefined{chapters@#1}%
    {\PackageError{chapters}{Sheet #1 is not listed in the index}%
      {List it in 20261125T012000--grh-workbook__guide.org\MessageBreak
       and run scripts/make-chapters.py.}}%
    {\edef\thechap{\@nameuse{chapters@#1}}}}
\makeatother
\chapnumber{20261125T014500}
\renewcommand{\thesection}{\thechap.\arabic{section}}

\theoremstyle{definition}
\newtheorem{problem}{Problem}[section]
\theoremstyle{plain}
\newtheorem{theorem}{Theorem}[section]

\newcommand{\Z}{\mathbb{Z}}
\newcommand{\Q}{\mathbb{Q}}
\newcommand{\F}{\mathbb{F}}
\newcommand{\Zp}{\Z_p}
\newcommand{\Oc}{\mathcal{O}}
\newcommand{\Gal}{\operatorname{Gal}}
\newcommand{\Hom}{\operatorname{Hom}}
\newcommand{\Ind}{\operatorname{Ind}}
\newcommand{\Frob}{\mathrm{Frob}}
\newcommand{\tr}{\operatorname{tr}}
\newcommand{\ad}{\operatorname{ad}}
\newcommand{\num}{\operatorname{num}}
\newcommand{\Sha}{\mathrm{Sha}}
\newcommand{\GQ}{G_{\Q,p}}
\newcommand{\mat}[4]{\bigl(\begin{smallmatrix} #1 & #2 \\ #3 & #4 \end{smallmatrix}\bigr)}

\setlength{\emergencystretch}{2em}
\title{Eisenstein congruences and Selmer classes}
\author{}
\date{November 25, 2026}

\begin{document}
\maketitle
\refstepcounter{chap}\label{chap}

\paragraph{Prerequisites.}
Modular forms of level \(1\), Eisenstein series, \(\Delta\) and the Hecke operators \cite[Ch.~VII]{SeA}; Serre writes the weight as \(2k\) and numbers the Bernoulli numbers differently.  Complete discrete valuation rings and Hensel's lemma.  The Chebotarev density theorem.  From Chapter~\ref{ds-chap}, \emph{Denominators and the Stickelberger element}: the eigenspaces \(A^{(i)}\) of Section~\ref{ds-sec:stickelberger} and Theorem~\ref{ds-thm:HR}.  From Chapter~\ref{tt-chap}, \emph{Odd Tate twists}, with \(\ell = p\): the groups \(\GQ\), \(G_p\), \(\Sha_p(\Q, e)\), Problems~\ref{tt-pr:mod-l} and~\ref{tt-pr:induced}, Theorem~\ref{tt-thm:sha-class}, and the Euler characteristic formula stated before Problem~\ref{tt-pr:sha2}.  From Chapter~\ref{cg-chap}, \emph{Congruence modules}: Problem~\ref{cg-pr:691}.  Galois representations attached to modular forms are stated in Section~\ref{sec:ribet} with references.

Chapter~\ref{tt-chap} computed the length of \(\Sha_p(\Q, e)\) for positive odd \(e\) from the theorem of Mazur and Wiles, which was used as a black box.  How does one produce classes in \(\Sha\), and how does one bound them?  Ribet's answer to the first question \cite{Ri} is that a cusp form congruent to an Eisenstein series gives a Galois representation that is reducible modulo \(p\), and a good choice of lattice turns it into an unramified class.  The subgroup of \(\Sha_p(\Q,k-1)\) killed by \(p\) has a deformation-theoretic interpretation as a space of tangent vectors at the reducible Eisenstein representation.  We work with level \(1\), and the example to keep in mind is \(k = 12\), \(p = 691\).  Throughout, \(p\) is an odd prime, \(\chi\) is the \(p\)-adic cyclotomic character and \(\omega\) is its reduction modulo \(p\).

\section{The congruence}\label{sec:congruence}

For even \(k \ge 4\) let \(\sigma_{k-1}(n) := \sum_{d \mid n} d^{k-1}\), and normalize the Eisenstein series as
\[
  E_k := 1 - \frac{2k}{B_k} \sum_{n \ge 1} \sigma_{k-1}(n)\, q^n ,
\]
so that \(E_k\) has rational coefficients and constant term \(1\) \cite[Ch.~VII, \S4]{SeA}.  Let \(\Delta = q\prod_{n\ge1}(1-q^n)^{24} = \sum \tau(n) q^n\).  Since \(B_{12} = -691/2730\), the least common denominator of the coefficients of \(E_{12}\) is \(691\).  By Problem~\ref{cg-pr:691} of Chapter~\ref{cg-chap}, \(691\,(E_4^3 - E_{12}) = 432000\,\Delta\).

\begin{problem}\label{pr:691}
\(\tau(n) \equiv \sigma_{11}(n) \pmod{691}\) for all \(n \ge 1\).
\end{problem}

The same happens whenever \(p\) divides \(B_k/k\).

\begin{problem}\label{pr:cuspform}
Let \(k \ge 4\) be even, and let \(p\) divide the numerator of \(B_k/k\).  There is a cusp form \(g\) of weight \(k\) and level \(1\) with coefficients in \(\Z_{(p)}\) such that \(a_n(g) \equiv \sigma_{k-1}(n) \pmod p\) for all \(n \ge 1\).
\end{problem}

\begin{problem}\label{pr:eigen-modp}
With \(g\) as in Problem~\ref{pr:cuspform}, for every prime \(\ell\) and every \(n \ge 1\),
\[
  a_n(T_\ell\, g) \equiv (1 + \ell^{k-1})\, a_n(g) \pmod p .
\]
\end{problem}

The cusp forms with coefficients in \(\Z_{(p)}\) form a lattice stable under the Hecke operators.  By the lifting lemma of Deligne and Serre \cite[Lemme~6.11]{DeS}, Problem~\ref{pr:eigen-modp} gives a normalized eigenform \(f\) of weight \(k\) and level \(1\), and a prime \(\lambda \mid p\) of the field generated by its coefficients, such that
\begin{equation}\label{eq:congruence}
  a_\ell(f) \equiv 1 + \ell^{k-1} \pmod \lambda \qquad \text{for every prime } \ell .
\end{equation}
In particular \(a_p(f) \equiv 1\): \(f\) is \emph{ordinary} at \(\lambda\).  For \(k = 12\) we may take \(f = \Delta\), and \(\tau(691) \equiv 1 \pmod{691}\).

\section{Lattices}\label{sec:lattices}

A congruence between \(f\) and an Eisenstein series says that a Galois representation is reducible modulo \(\lambda\).  Which reduction one sees depends on a lattice.  Let \(K\) be a finite extension of \(\mathbb{Q}_p\) with ring of integers \(\Oc\), uniformizer \(\pi\), valuation \(v\) with \(v(\pi) = 1\), and residue field \(\F\).  Let \(G\) be a profinite group and \(\rho \colon G \to \mathrm{GL}(V)\) a continuous irreducible representation on a \(2\)-dimensional \(K\)-space.  A \emph{lattice} is a free \(\Oc\)-submodule of rank \(2\) that spans \(V\).

\begin{problem}\label{pr:stable}
\(G\)-stable lattices exist.
\end{problem}

Suppose that for one, hence (by the Brauer--Nesbitt theorem \cite[(30.16)]{CR}) every, \(G\)-stable lattice \(L\), the semisimplification of \(L/\pi L\) is \(\chi_1 \oplus \chi_2\) with characters \(\chi_1 \ne \chi_2 \colon G \to \F^\times\).

\begin{problem}\label{pr:diag}
Choose \(\sigma_0 \in G\) with \(\chi_1(\sigma_0) \ne \chi_2(\sigma_0)\).  Then \(\rho(\sigma_0)\) has eigenvalues \(\lambda_1, \lambda_2 \in \Oc\) with \(\lambda_i \equiv \chi_i(\sigma_0) \pmod \pi\), and every \(G\)-stable lattice \(L\) has a basis \(e_1, e_2\) with \(\rho(\sigma_0) e_i = \lambda_i e_i\).
\end{problem}

Fix such \(L\), \(e_1\), \(e_2\), write \(\rho(\sigma) = \mat{a(\sigma)}{b(\sigma)}{c(\sigma)}{d(\sigma)}\) in this basis, and put \(m := \min_\sigma v(b(\sigma))\), \(n := \min_\sigma v(c(\sigma))\).  By irreducibility \(b\) and \(c\) are not identically zero, so \(m, n\) are finite.  For \(r \in \Z\) let \(L_r := \pi^r \Oc e_1 \oplus \Oc e_2\).

\begin{problem}\label{pr:segment}
\(L_r\) is \(G\)-stable if and only if \(-n \le r \le m\), and every \(G\)-stable lattice is homothetic to exactly one \(L_r\): the stable lattices form a segment of length \(m+n\) in the tree of \(\mathrm{PGL}_2(K)\).
\end{problem}

\begin{problem}\label{pr:ends}
\(m + n \ge 1\).  On \(L_m/\pi L_m\) the group \(G\) acts by \(\mat{\chi_1}{*}{0}{\chi_2}\), a nonsplit extension; on \(L_{-n}/\pi L_{-n}\) it acts by a nonsplit \(\mat{\chi_1}{0}{*}{\chi_2}\), with sub \(\chi_2\) and quotient \(\chi_1\); and for \(-n < r < m\) it acts by \(\chi_1 \oplus \chi_2\).
\end{problem}

Problems~\ref{pr:stable}--\ref{pr:ends} prove the following.  It is due to Ribet \cite[Prop.~2.1]{Ri}.

\begin{theorem}[Ribet's lemma]\label{thm:ribet-lemma}
There is a \(G\)-stable lattice on whose reduction \(G\) acts by a nonsplit extension \(\mat{\chi_1}{*}{0}{\chi_2}\).
\end{theorem}

\section{The unramified class}\label{sec:ribet}

We state three inputs.  Let \(f\) be a normalized cuspidal Hecke eigenform of weight \(k\ge2\) and level \(1\), let \(\lambda \mid p\) be a prime of the field of its coefficients, and let \(K\) be the completion at \(\lambda\).
\begin{itemize}
\item[(D)] There is a continuous \(\rho_f \colon \GQ \to \mathrm{GL}_2(K)\) with \(\tr \rho_f(\Frob_\ell) = a_\ell(f)\) and \(\det \rho_f = \chi^{k-1}\), for every prime \(\ell \ne p\), where \(\Frob_\ell\) denotes arithmetic Frobenius \cite{De}.
\item[(I)] \(\rho_f\) is irreducible \cite[\S2]{Ri2}.
\item[(O)] If \(a_p(f)\) is a \(\lambda\)-adic unit, then \(\rho_f|_{G_p}\) stabilizes a line on which the inertia group \(I_p\) acts by \(\chi^{k-1}\), and \(G_p\) acts on the quotient by an unramified character \cite[Thm.~2.1.4]{Wi}.
\end{itemize}
Chapter~\ref{gr-chap}, \emph{Black box: Galois representations of eigenforms}, states the analogue of (D) as a black box, over CM fields; (I) and (O) are not part of it.  Its Section~\ref{gr-sec:es} shows in the simplest case, the eigenform \(f\) of weight \(2\) and level \(11\), where \(\rho_f\) comes from: it is the rational \(p\)-adic Tate module of \(J_0(11)\), and the Eichler--Shimura relation gives the analogue of (D).

From now on \(k\) is even with \(4 \le k \le p - 3\), \(p\) divides \(B_k\) (equivalently, since \(k < p\), the numerator of \(B_k/k\)), and \(f\), \(\lambda\) are as in \eqref{eq:congruence}.

\begin{problem}\label{pr:ss}
The semisimplification of \(\rho_f\) modulo \(\lambda\) is \(1 \oplus \omega^{k-1}\), and \(1 \ne \omega^{k-1}\) on \(G_p\), and even on \(I_p\).
\end{problem}

Apply Theorem~\ref{thm:ribet-lemma} with \(\chi_1 = 1\) and \(\chi_2 = \omega^{k-1}\), and let \(\bar\rho = \mat{1}{\bar b}{0}{\omega^{k-1}}\) be the resulting reduction.

\begin{problem}\label{pr:class}
\(\kappa(\sigma) := \bar b(\sigma)\,\omega^{1-k}(\sigma)\) is a cocycle on \(\GQ\) with values in \(\F(\omega^{1-k})\), and its class is not zero.
\end{problem}

\begin{problem}\label{pr:unram}
The class of \(\kappa\) vanishes on \(G_p\).
\end{problem}

Since \(H^1(\GQ, \F(\omega^{1-k})) = H^1(\GQ, \Z/p(1-k)) \otimes_{\F_p} \F\), and similarly for \(G_p\), Problems~\ref{pr:ss}--\ref{pr:unram}, Problem~\ref{tt-pr:mod-l} and Theorem~\ref{tt-thm:sha-class} prove the following.

\begin{theorem}[Ribet]\label{thm:ribet}
Let \(k\) be even with \(4 \le k \le p-3\), and suppose that \(p \mid B_k\).  Then \(\Sha_p(\Q, k-1) \ne 0\), and \(A^{(p-k)} \ne 0\).
\end{theorem}

This is the converse to Herbrand's theorem \cite{Ri}; see also \cite[Thm.~6.18]{Wa}.  For \(k = 12\) and \(p = 691\) the input is \(\Delta\).  There is an unramified cyclic extension of \(\Q(\zeta_{691})\) of degree \(691\), Galois over \(\Q\), with conjugation action on its relative Galois group through \(\omega^{-11} = \omega^{679}\).  This proves nonvanishing of \(\Sha_{691}(\Q,11)\).  Congruences with Eisenstein series are measured by the Eisenstein ideal of the Hecke algebra, which Mazur introduced in weight \(2\) and prime level \cite{Ma}.  In this range Theorem~\ref{tt-thm:MW}, a consequence of the main conjecture of Iwasawa theory, proved by Mazur and Wiles \cite{MW}, gives \(\operatorname{length}_{\Zp}\Sha_p(\Q,k-1)=v_p(B_k/k)\).

\section{Tangent vectors at the Eisenstein point}\label{sec:tangent}

Ribet goes from a modular form to a class.  Here a congruence produces a nonzero class.  How does a class in \(\Sha_p(\Q,k-1)[p]\) give a first-order deformation?  Let \(\rho_0 := 1 \oplus \omega^{k-1} \colon \GQ \to \mathrm{GL}_2(\F_p)\), the reduction of the representation \(1 \oplus \chi^{k-1}\) of the Eisenstein series \(E_k\); we call it the \emph{Eisenstein point}.  Let \(\F_p[\varepsilon]\) be the ring with \(\varepsilon^2 = 0\).  A \emph{lift} of \(\rho_0\) is a continuous homomorphism \(\rho_\varepsilon \colon \GQ \to \mathrm{GL}_2(\F_p[\varepsilon])\) that reduces to \(\rho_0\) modulo \(\varepsilon\).  Let \(\ad \rho_0\) be the space \(M_2(\F_p)\) with \(\sigma\) acting by conjugation by \(\rho_0(\sigma)\).  See \cite[\S1]{Mz} for the general theory.

\begin{problem}\label{pr:tangent}
Writing \(\rho_\varepsilon = (1 + \varepsilon M)\rho_0\) identifies the lifts of \(\rho_0\), up to conjugation by matrices \(1 + \varepsilon X\), with \(H^1(\GQ, \ad\rho_0)\).
\end{problem}

The torus \(T = \{\mathrm{diag}(1, t)\}\) commutes with \(\rho_0\), and acts on lifts by conjugation.

\begin{problem}\label{pr:weights}
\(\ad\rho_0 = \F_p E_{11} \oplus \F_p E_{22} \oplus \F_p(\omega^{1-k}) E_{12} \oplus \F_p(\omega^{k-1}) E_{21}\), where \(E_{ij}\) are the matrix units, and \(T\) acts on these four summands, and on their \(H^1\), with weights \(0\), \(0\), \(-1\), \(+1\).
\end{problem}

A cocycle \(\kappa\) with nonzero class in \(H^1(\GQ, \F_p(\omega^{1-k}))\) gives the nonsplit \(\rho_\kappa = \mat{1}{\kappa\omega^{k-1}}{0}{\omega^{k-1}}\).  For \(t \in \F_p^\times\), conjugating by \(\mathrm{diag}(t,1)\) replaces \(\kappa\) by \(t\kappa\), so all the \(\rho_{t\kappa}\) are isomorphic; at \(t = 0\) one reaches \(\rho_0\), and to first order at \(t = 0\) the family is the lift \((1 + \varepsilon\kappa E_{12})\rho_0\), of weight \(-1\).  This is the sense in which a class is a tangent vector at the Eisenstein point.

Which lifts are allowed?  By (O), \(\rho_f\) has a line on which \(I_p\) acts by \(\chi^{k-1}\) with unramified quotient.  Call a lift \(\rho_\varepsilon\) \emph{ordinary} if \(\F_p[\varepsilon]^2\) has a \(G_p\)-stable direct summand \(L\) of rank \(1\), reducing to \(\F_p e_2\), on which \(I_p\) acts by \(\omega^{k-1}\), with \(I_p\) acting trivially on the quotient.  Let \(t_{\mathrm{ord}} \subseteq H^1(\GQ, \ad\rho_0)\) be the set of classes of ordinary lifts, and \(t_{\mathrm{ord}}^{(j)}\) its part of weight \(j\).

\begin{problem}\label{pr:ordinary}
The lift with \(M = \mat{\alpha_1}{\beta}{\gamma}{\alpha_2}\) is ordinary if and only if \(\alpha_1\) and \(\alpha_2\) vanish on \(I_p\) and the class of \(\beta\) vanishes in \(H^1(G_p, \F_p(\omega^{1-k}))\).  Hence \(t_{\mathrm{ord}}^{(0)} = 0\), and \(t_{\mathrm{ord}}^{(+1)} = H^1(\GQ, \F_p(\omega^{k-1}))\).
\end{problem}

With Problem~\ref{tt-pr:mod-l} and Theorem~\ref{tt-thm:sha-class}, Problems~\ref{pr:tangent}--\ref{pr:ordinary} prove the following.

\begin{theorem}\label{thm:tangent}
\(t_{\mathrm{ord}} = t_{\mathrm{ord}}^{(-1)} \oplus t_{\mathrm{ord}}^{(+1)}\), and \(t_{\mathrm{ord}}^{(-1)} \cong \Sha_p(\Q, k-1)[p] \cong \Hom(A^{(p-k)}, \F_p)\).
\end{theorem}

\begin{problem}\label{pr:plus-side}
\(\dim t_{\mathrm{ord}}^{(+1)} \ge 1\).
\end{problem}

In (D) the Hecke eigenvalues at primes \(\ell\ne p\) are traces of arithmetic Frobenius.

\begin{problem}\label{pr:trace-blind}
Every ordinary lift of \(\rho_0\) to \(\F_p[\varepsilon]\) has trace \(1 + \omega^{k-1}\).
\end{problem}

The traces see \(\beta\) and \(\gamma\) only through their products.  Return to \(\rho_f\) with the basis \(e_1, e_2\) of Problem~\ref{pr:diag}, with \(\chi_1 = 1\) and \(\chi_2 = \omega^{k-1}\), and let \(I \subseteq \Oc\) be the ideal generated by the products \(b(\sigma)c(\tau)\), \(\sigma, \tau \in \GQ\).  Following Bellaïche and Chenevier \cite[Ch.~1]{BC}, we call it the \emph{reducibility ideal}.

\begin{problem}\label{pr:red-ideal}
\(I = \pi^{m+n}\Oc\), and \(a\) and \(d\) are characters of \(\GQ\) modulo \(I\).
\end{problem}

\begin{problem}\label{pr:trace-cong}
\(a \equiv 1\) and \(d \equiv \chi^{k-1}\) modulo \(I\).  Hence \(v(a_\ell(f) - 1 - \ell^{k-1}) \ge m + n\) for every prime \(\ell \ne p\).
\end{problem}

\begin{problem}\label{pr:691-segment}
For \(f = \Delta\) and \(p = 691\), \(m + n = 1\): up to homothety \(\rho_\Delta\) has exactly two stable lattices.  Moreover \(v_{691}(\tau(\ell) - 1 - \ell^{11}) = 1\) for every prime \(\ell \le 37\).
\end{problem}

The valuations of the trace differences bound the length \(m+n\) of the segment of stable lattices from above.  To seek a converse bound by modular methods one can compare an integral deformation ring with a Hecke algebra.  For modularity theorems under ordinary residually reducible hypotheses over \(\Q\), see \cite{SW}.  For comparison with Problem~\ref{pr:691-segment}, Theorem~\ref{tt-thm:MW} gives \(\operatorname{length}_{\Z_{691}}\Sha_{691}(\Q,11)=1\).

\paragraph{Further reading.}
Chapter~\ref{tf-chap}, \emph{Twists with few congruences}, treats simultaneous twisting, and Chapter~\ref{oc-chap}, \emph{Odd classes by Gauss sums}, treats the bound for odd classes.  The rewritten Lean proof \cite{GL} gives the projection denominator bound in \texttt{Selmer.denominatorLe\_clearingBound}, in \texttt{grh/src/New/Congruences/FirstOrder/DenominatorBound.lean}, and the simultaneous twisting theorem in \texttt{Selmer.exists\_mFoldTwist\_hasDivisibleSelmerLengthLe}, in \texttt{grh/src/New/Congruences/Selmer/TwistingTheorem.lean}; here and below Lean names omit the namespace prefix \texttt{New.}

\section{From the CM field to characters}\label{sec:shapiro}

The two results above use Galois representations over the CM field \texttt{Data.E} of Chapter~\ref{c32-chap}, \emph{The data}.  This field is abelian over \(\Q\) by the cyclotomic embedding \texttt{Parameters.toCycloField} \cite{GL}.  In this section \(E\) denotes an arbitrary abelian number field.  For the abelian fields and coefficient ring below, Shapiro's lemma and the character idempotents decompose a class over \(E\) into a family of character-twisted classes over \(\Q\).  Problem~\ref{tt-pr:induced} treats the cyclotomic case.

Let \(E\) be an abelian number field with \(\Gamma := \Gal(E/\Q)\) of order prime to \(p\).  Let \(\Oc := \Zp[\mu_{\lvert\Gamma\rvert}]\), let \(W := W_e \otimes_{\Zp} \Oc\) with \(W_e = \mathbb{Q}_p/\Zp(-e)\) as in Chapter~\ref{tt-chap}, and let \(G_E \subseteq G_\Q\) be the absolute Galois groups.  Let \(\Sha_p(E, e)_\Oc\) be the group of classes in \(H^1(G_E, W)\) that vanish on the inertia group \(I_w\) of every place \(w \nmid p\) and on the decomposition group \(D_w\) of every \(w \mid p\).  Since \(\Oc\) is free of finite rank over \(\Zp\), continuous cohomology and the restriction maps commute with \(\otimes_{\Zp} \Oc\), so \(\Sha_p(E, e)_\Oc = \Sha_p(E, e) \otimes_{\Zp} \Oc\), with \(\Sha_p(E, e)\) as in Chapter~\ref{tt-chap}.  For a character \(\psi \colon \Gamma \to \Oc^\times\), let \(\Sha_p(\Q, e; \psi)_\Oc\) be defined in the same way from \(H^1(G_\Q, W \otimes \psi^{-1})\).  Let \(\Gamma\) act on \(\Oc[\Gamma]\) by left multiplication, so that \(G_\Q\) acts on it through \(\Gamma\), and let \(\gamma \in \Gamma\) act on \(H^1(G_E, W)\) by \((\gamma c)(h) := \tilde\gamma\, c(\tilde\gamma^{-1} h \tilde\gamma)\), for any lift \(\tilde\gamma \in G_\Q\) of \(\gamma\).

\begin{problem}\label{pr:shapiro}
\(H^1(G_E, W) \cong H^1(G_\Q, \Oc[\Gamma] \otimes_\Oc W)\), and \(\Oc[\Gamma] \otimes_\Oc W \cong \bigoplus_\psi W \otimes \psi\) as \(G_\Q\)-modules, the sum over all characters \(\psi\) of \(\Gamma\).
\end{problem}

\begin{problem}\label{pr:mackey}
Let \(v\) be a prime of \(\Q\), and \(D = D_v\) or \(I_v\).  A class of \(H^1(G_E, W)\) vanishes on \(D_w\), respectively on \(I_w\), for every place \(w \mid v\) of \(E\) if and only if its image under Problem~\ref{pr:shapiro} vanishes on \(D\).
\end{problem}

\begin{problem}\label{pr:sha-decomp}
\(\Sha_p(E, e)_\Oc \cong \bigoplus_\psi \Sha_p(\Q, e; \psi)_\Oc\).  The summand \(\Sha_p(\Q, e; \psi)_\Oc\), which comes from the summand \(W \otimes \psi^{-1}\) of Problem~\ref{pr:shapiro}, is the part on which \(\Gamma\) acts through \(\psi\).
\end{problem}

Every class killed by \(p\) in the \(\psi\)-part of \(\Sha_p(E, e)_\Oc\) lifts to \(H^1(G_\Q, \F(\omega^{-e}\bar\psi^{-1}))\), with \(\F\) the residue field of \(\Oc\).  As in \S\ref{sec:tangent}, a cocycle representing this lift gives an upper-triangular deformation of \(1\oplus\omega^e\bar\psi\) over \(\F[\varepsilon]\), with \(\varepsilon^2=0\).  The deformation allows ramification at \(p\) and at primes dividing the conductor of \(\psi\).  When the global and local residual invariants vanish, the coefficient maps are injective there and the strict local condition passes to the lift, as in Problem~\ref{tt-pr:mod-l}.  For \(E=\Q\) and positive odd \(e\), Theorem~\ref{tt-thm:MW} gives the length \(v_p(\num(B_{e+1}/(e+1)))\).  For \(E\subseteq\Q(\zeta_p)\), write the reduction of \(\psi\) as \(\omega^j\) after inflation to \(\Gal(\Q(\zeta_p)/\Q)\).  If \(i\equiv-e-j\not\equiv0\pmod{p-1}\), Problem~\ref{tt-pr:eta-rank} identifies the subgroup killed by \(p\) in this summand with \(\Hom(A^{(i)},\F_p)\).  For odd \(3\le i\le p-2\), Theorem~\ref{ds-thm:HR} gives the order of \(A^{(i)}\) and its criterion for nonvanishing.  The case \(i=1\) is discussed after Problem~\ref{ds-pr:gcd-c}.

\section*{Hints}

\subsection*{\S\ref{sec:congruence}}
\begin{description}
\item[\ref{pr:691}] Reduce the identity modulo \(691\), noting \(432000 + 65520 = 720 \cdot 691\).
\item[\ref{pr:cuspform}] The forms \(E_4^a E_6^b\).
\item[\ref{pr:eigen-modp}] Compare the powers of \(\ell\) dividing the indices in the Hecke coefficient formula.
\end{description}

\subsection*{\S\ref{sec:lattices}}
\begin{description}
\item[\ref{pr:stable}] Consider the orbit of a lattice under the compact group \(\rho(G)\) in the discrete space of lattices.
\item[\ref{pr:diag}] Hensel's lemma for the characteristic polynomial.  The projector \((\rho(\sigma_0) - \lambda_2)/(\lambda_1 - \lambda_2)\) has entries in \(\Oc\).
\item[\ref{pr:segment}] By Problem~\ref{pr:diag}, a stable lattice is the sum of its intersections with \(Ke_1\) and \(Ke_2\).  Write \(\rho\) in the basis \(\pi^r e_1, e_2\).
\item[\ref{pr:ends}] The only lines stable under \(\bar\rho(\sigma_0)\) are \(\F e_1\) and \(\F e_2\).  A representation with semisimplification \(\chi_1 \oplus \chi_2\) has a stable line.
\end{description}

\subsection*{\S\ref{sec:ribet}}
\begin{description}
\item[\ref{pr:ss}] Chebotarev, Brauer--Nesbitt, and \eqref{eq:congruence}.  \(\omega\) maps \(I_p\) onto \(\F_p^\times\).
\item[\ref{pr:class}] Multiply out \(\bar\rho(\sigma\tau)\).  A coboundary can be removed by conjugating with a unipotent upper triangular matrix.
\item[\ref{pr:unram}] Reduce the line of (O) inside the lattice of Theorem~\ref{thm:ribet-lemma}, and compare it with \(\F e_1\) under \(I_p\).
\end{description}

\subsection*{\S\ref{sec:tangent}}
\begin{description}
\item[\ref{pr:tangent}] Expand \(\rho_\varepsilon(\sigma\tau) = \rho_\varepsilon(\sigma)\rho_\varepsilon(\tau)\).
\item[\ref{pr:weights}] Conjugate each \(E_{ij}\) by \(\rho_0(\sigma)\) and by \(\mathrm{diag}(1,t)\).
\item[\ref{pr:ordinary}] A rank-one summand reducing to \(\F_p e_2\) is spanned by \(e_2 + \varepsilon x e_1\).  For the weight \(0\) part, \(\Q\) has no nontrivial extension unramified at every prime.
\item[\ref{pr:plus-side}] The Euler characteristic formula of Chapter~\ref{tt-chap}.
\item[\ref{pr:trace-blind}] Problem~\ref{pr:ordinary}.
\item[\ref{pr:red-ideal}] The upper left entry of \(\rho(\sigma\tau)\).
\item[\ref{pr:trace-cong}] In the basis of \(L_m\), \(\rho\) is upper triangular modulo \(I\).  Compare the line \((\Oc/I)e_1\) with the reduction of the line of (O) under \(I_p\).  A character of \(\GQ\) unramified at \(p\) is trivial.  Then use the determinant.
\item[\ref{pr:691-segment}] Problems~\ref{pr:ends} and~\ref{pr:trace-cong} at \(\ell = 2\), where \(\tau(2) = -24\).  The last statement is a finite computation with the product for \(\Delta\).
\end{description}

\subsection*{\S\ref{sec:shapiro}}
\begin{description}
\item[\ref{pr:shapiro}] Shapiro's lemma.  The elements \(\lvert\Gamma\rvert^{-1}\sum_\gamma \psi(\gamma)^{-1}\gamma\) are orthogonal idempotents.
\item[\ref{pr:mackey}] Mackey's formula.  The double cosets \(D_v \backslash G_\Q / G_E\) correspond to the places of \(E\) above \(v\); for \(I_v\) each place \(w\) occurs \(f(w/v)\) times, its residue degree, through conjugates of \(I_w\).
\item[\ref{pr:sha-decomp}] The decomposition of Problem~\ref{pr:shapiro} is \(G_\Q\)-stable, hence stable under every \(D\).  Under Problem~\ref{pr:shapiro}, \(\gamma\) acts on \(\Oc[\Gamma] \otimes_\Oc W\) by right multiplication by \(\gamma^{-1}\) on \(\Oc[\Gamma]\).
\end{description}

\begin{thebibliography}{DeS}
\bibitem[BC]{BC} J.~Bella\"iche and G.~Chenevier, \emph{Families of Galois representations and Selmer groups}, Ast\'erisque \textbf{324} (2009).
\bibitem[CR]{CR} C.~W.~Curtis and I.~Reiner, \emph{Representation Theory of Finite Groups and Associative Algebras}, Interscience, 1962.
\bibitem[De]{De} P.~Deligne, \emph{Formes modulaires et repr\'esentations \(l\)-adiques}, S\'eminaire Bourbaki 1968/69, exp.~355, Lecture Notes in Math.\ 179, Springer, 1971, 139--172.
\bibitem[DeS]{DeS} P.~Deligne and J.-P.~Serre, \emph{Formes modulaires de poids \(1\)}, Ann.\ Sci.\ \'Ec.\ Norm.\ Sup.\ (4) \textbf{7} (1974), 507--530.
\bibitem[GL]{GL} L.~Sernau, \emph{GRH for Dirichlet \(L\)-functions, in Lean~4}, a shorter rewriting of the Lean proof, text by Claude, \url{https://github.com/lsernau/grh-lean}, commit \texttt{72fead1}, readable edition \texttt{grh/}.
\bibitem[Ma]{Ma} B.~Mazur, \emph{Modular curves and the Eisenstein ideal}, Publ.\ Math.\ IH\'ES \textbf{47} (1977), 33--186.
\bibitem[MW]{MW} B.~Mazur and A.~Wiles, \emph{Class fields of abelian extensions of \(\mathbf{Q}\)}, Invent.\ Math.\ \textbf{76} (1984), 179--330.
\bibitem[Mz]{Mz} B.~Mazur, \emph{Deforming Galois representations}, in: Galois Groups over \(\mathbf{Q}\), Math.\ Sci.\ Res.\ Inst.\ Publ.\ 16, Springer, 1989, 385--437.
\bibitem[Ri]{Ri} K.~A.~Ribet, \emph{A modular construction of unramified \(p\)-extensions of \(\mathbf{Q}(\mu_p)\)}, Invent.\ Math.\ \textbf{34} (1976), 151--162.
\bibitem[Ri2]{Ri2} K.~A.~Ribet, \emph{Galois representations attached to eigenforms with Nebentypus}, in: Modular Functions of One Variable V, Lecture Notes in Math.\ 601, Springer, 1977, 17--52.
\bibitem[SeA]{SeA} J.-P.~Serre, \emph{A Course in Arithmetic}, Grad.\ Texts in Math.\ 7, Springer, 1973.
\bibitem[SW]{SW} C.~M.~Skinner and A.~J.~Wiles, \emph{Residually reducible representations and modular forms}, Publ.\ Math.\ IH\'ES \textbf{89} (1999), 5--126.
\bibitem[Wa]{Wa} L.~C.~Washington, \emph{Introduction to Cyclotomic Fields}, 2nd ed., Grad.\ Texts in Math.\ 83, Springer, 1997.
\bibitem[Wi]{Wi} A.~Wiles, \emph{On ordinary \(\lambda\)-adic representations associated to modular forms}, Invent.\ Math.\ \textbf{94} (1988), 529--573.
\end{thebibliography}

\end{document}
